\adj{} computes the sensitivity of a performance measure to model inputs by
solving an adjoint problem alongside a \mf{} simulation \citep{hayek2025}. The
approach is non-intrusive: the forward model runs unmodified through the
\mf{} application programming interface, and the quantities the adjoint needs
are read from the memory manager as the simulation proceeds.

This document records the technical detail behind extensions made after the
original publication. It is supplemental: it assumes the formulation in
\citet{hayek2025} and describes only what has been added or corrected. Each
chapter is self-contained and can be read on its own.

\section{Scope}

Chapter~\ref{ch:instantaneous} describes the instantaneous performance measure,
a third form alongside the direct and residual forms, which averages a quantity
over the time steps at which it is observed rather than accumulating it.

Chapter~\ref{ch:storage} describes how the storage terms couple one time step to
the next in the adjoint, and how the specific-storage and specific-yield
contributions are selected to match what \mf{} does.

Chapters~\ref{ch:lake} and \ref{ch:sfr} describe the treatment of two advanced
packages whose dependent variables are solved outside the groundwater flow
matrix. Both use the same device, a bordered matrix, and the two chapters share
the notation established in Chapter~\ref{ch:lake}.

\section{Notation}

The groundwater flow equations for one time step are written

\begin{equation}
	\label{eqn:forward}
	\mathbf{A} \mathbf{h} = \mathbf{b},
\end{equation}

\noindent where $\mathbf{A}$ is the flow matrix, $\mathbf{h}$ is the vector of
cell heads, and $\mathbf{b}$ is the right-hand side. For a performance measure
$f$ the adjoint state $\lam$ satisfies

\begin{equation}
	\label{eqn:adjoint}
	\mathbf{A}^{T} \lam = \frac{\partial f}{\partial \mathbf{h}},
\end{equation}

\noindent and the sensitivity of $f$ to a parameter $p$ that enters through the
residual $\mathbf{r} = \mathbf{A}\mathbf{h} - \mathbf{b}$ is

\begin{equation}
	\label{eqn:sensitivity}
	\frac{\partial f}{\partial p} = -\lam^{T} \frac{\partial \mathbf{r}}
	{\partial p}.
\end{equation}

A well withdraws at a specified volumetric rate, so
$\partial \mathbf{r} / \partial q$ is a unit vector at the pumped cell and the
sensitivity of $f$ to that rate is the adjoint state itself. That identity is
used throughout as a check, because it means the reported well sensitivity is
$\lam$ without post-processing.

Superscripts $t$ and $\told$ denote the current and previous time step. The
saturated fraction of a cell is $S_F$, its horizontal area is $A$, and its
thickness is $\Delta z$.
